%! Function Fundamentals — Unit 1, Lesson 1.1 — Homework (ANSWER KEY)
% THE LESSON'S GRADED PRACTICE and the last component in the packet — exactly TWO pages, started
% in the Homework Launch (the last 3 minutes of the period) and finished at home. Every context
% here is new: the notes used a rideshare fare, a vending machine, and an unlabeled graph; the
% homework uses a phone plan, a temperature log, a new graph, and a new formula.
%   Page 1 — Part A (evaluating, 1–3) and Part B (a table read both ways, 4).
%   Page 2 — Part C (a graph read both ways, 5–7) and Part D (domain and range, 8–10).
% The diagnostic item is 4(c): solve T(h) = 70 has NO answer — the "zero" in "zero, one, or
% many" that the notes only named. Problem 10 asks for the lesson's headline in writing.
%
% PAGE LOCKSTEP with the other half of the pair: work blocks are byte-identical, prose answers
% are \writespace{H}{…} (fixed height both ways), and each \blank sits at the END of a line so
% its \ans cannot rewrap anything. The explicit \newpage fixes the page break in both files.
\documentclass[10pt]{article}
\usepackage{precalculus-article}
\usepackage{precalculus-key}
\newcommand{\TallMath}[1]{$\displaystyle #1\rule[-1.4em]{0pt}{3.2em}$}

\begin{document}

\pageheader{Unit 1, Lesson 1.1}{Homework --- Answer Key}

\vspace{0.12in}

\boxguard[8]
\begin{remindbox}
\small
\textbf{This is your graded homework.} Start it in class; finish it on your own by the due
date on your cover sheet. Show your work: fill the slot with parentheses before you simplify,
and read every question's \emph{direction} --- evaluate or solve --- before you compute.
\end{remindbox}

\vspace{0.12in}

\begin{headlinebox}{lilac}
\textbf{\color{plum}Part A --- Evaluating from a formula}
\end{headlinebox}

\begin{enumerate}[leftmargin=1.9em, itemsep=10pt]

\item A phone plan charges a monthly fee plus a charge for each gigabyte of data:
      $C(g) = 15 + 8g$, where $C(g)$ is the monthly cost in dollars for $g$ gigabytes.

      \textbf{(a)} Write in words what $C(3)$ means --- no computing yet.
      \writespace{1.1cm}{The monthly cost, in dollars, of a month that uses 3 gigabytes of data.}

      \textbf{(b)} Find $C(3)$.

\begin{work}
  C(3) &= 15 + 8(3) \\
       &= 15 + 24 = 39
\end{work}

      $C(3) = $ \ans{\$39}

\item \begin{minipage}[t]{0.44\linewidth}
      Let $h(x) = 2x^2 - 5x$. Find $h(-3)$.

\begin{work}
  h(-3) &= 2(-3)^2 - 5(-3) \\
        &= 2(9) + 15 \\
        &= 33
\end{work}

      $h(-3) = $ \ans{33}
      \end{minipage}\hfill
      \begin{minipage}[t]{0.50\linewidth}
      \textbf{3.}\ Let $q(t) = -t^2 + 6$. Find $q(-2)$ and $q(0)$. \textit{Careful: square
      first, then apply the negative.}

\begin{work}
  q(-2) &= -(-2)^2 + 6 \\
        &= -4 + 6 = 2
\end{work}

      $q(-2) = $ \ans{2} \qquad $q(0) = $ \ans{6}
      \end{minipage}
\setcounter{enumi}{3}

\end{enumerate}

\vspace{0.14in}

\begin{headlinebox}{lilac}
\textbf{\color{plum}Part B --- A table, read in both directions}
\end{headlinebox}

\begin{enumerate}[leftmargin=1.9em, itemsep=8pt, start=4]

\item A hiker logs the air temperature $T(h)$, in $^\circ$F, $h$ hours after sunrise.

      \smallskip
      \begin{center}
      \renewcommand{\arraystretch}{1.35}
      \begin{tabular}{|c|c|c|c|c|c|c|}
      \hline
      $h$ (hours) & 0 & 1 & 2 & 3 & 4 & 5 \\
      \hline
      $T(h)$ ($^\circ$F) & 52 & 58 & 63 & 63 & 60 & 55 \\
      \hline
      \end{tabular}
      \end{center}

      \textbf{(a)} Find $T(4)$. \quad $T(4) = $ \ans{60 $^\circ$F} \hspace{3em}
      \textbf{(b)} Solve $T(h) = 63$. \quad $h = $ \ans{2 and 3}

      \textbf{(c)} Solve $T(h) = 70$. How many answers are there, and why?
      \writespace{1.1cm}{None --- no hour in the table has an output of 70. A solve question can have zero answers.}

      \textbf{(d)} Which of (a), (b), (c) was an \emph{evaluate} question? How can you tell?
      \writespace{1.1cm}{(a) --- it gave the input, 4, and asked for the output. (b) and (c) gave an output.}

\end{enumerate}

\newpage

\begin{headlinebox}{lilac}
\textbf{\color{plum}Part C --- A graph, read in both directions}
\end{headlinebox}

\begin{minipage}[t]{0.47\linewidth}
\vspace{0pt}
\centering
\begin{tikzpicture}[x=0.62cm, y=0.5cm, >=stealth]
  \draw[linegray, very thin, step=1] (-4,-2) grid (6,6);
  \draw[->, thick] (-4,0) -- (6.5,0) node[below right] {\small $x$};
  \draw[->, thick] (0,-2) -- (0,6.5) node[above] {\small $g(x)$};
  \foreach \x in {-3,-2,-1,1,2,3,4,5} \node[below, font=\scriptsize] at (\x,0) {\x};
  \foreach \y in {-1,1,2,3,4,5} \node[left, font=\scriptsize] at (0,\y) {\y};
  \draw[very thick, plum] (-3,3) -- (-1,-1) -- (2,5) -- (5,2);
  \fill[plum] (-3,3) circle (2.5pt);
  \fill[plum] (5,2) circle (2.5pt);
\end{tikzpicture}
\end{minipage}\hfill
\begin{minipage}[t]{0.5\linewidth}
\vspace{0pt}
\begin{enumerate}[leftmargin=1.9em, itemsep=10pt, start=5]
\item Find each value.

      $g(2) = $ \ans{5}

      $g(-1) = $ \ans{$-1$}

\item Solve each equation. Some have more than one answer.

      $g(x) = 3$: \quad $x = $ \ans{$-3$, 1, and 4}

      $g(x) = 1$: \quad $x = $ \ans{$-2$ and 0}

\item State the domain and range of $g$ as inequalities, then with brackets.

      Domain: \ans{$-3 \le x \le 5$, $[-3, 5]$}

      Range: \ans{$-1 \le y \le 5$, $[-1, 5]$}
\end{enumerate}
\end{minipage}

\vspace{0.16in}

\begin{headlinebox}{lilac}
\textbf{\color{plum}Part D --- Domain, range, and the two directions}
\end{headlinebox}

\begin{enumerate}[leftmargin=1.9em, itemsep=10pt, start=8]

\item A function is given as a list of points: $\{(-1, 4),\ (2, 0),\ (5, 4),\ (6, -3)\}$.

      Domain: \ans{$\{-1, 2, 5, 6\}$} \qquad Range: \ans{$\{-3, 0, 4\}$}

      Two points share the output 4. Is this still a function? Explain.
      \writespace{1.1cm}{Yes. Two inputs may share an output (many-to-one); no input has two outputs.}

\item Let $r(x) = \dfrac{12}{x + 4}$.

      \textbf{(a)} Find $r(2)$. \quad $r(2) = $ \ans{$12/6 = 2$}

      \textbf{(b)} Which input is forbidden, and why? Write the domain of $r$ in words.
      \writespace{1.4cm}{$x = -4$: the denominator would be 0, and division by zero is undefined. Domain: every real number except $-4$.}

\item A classmate says: ``$f$ is a function, so the equation $f(x) = 6$ can only have one
      answer.'' Use the words \emph{input} and \emph{output} to explain what is wrong.
      \writespace{1.9cm}{A function gives each INPUT exactly one output. Solving $f(x) = 6$ asks for every input whose output is 6 --- and several inputs can share one output, so there can be zero, one, or many answers.}

\end{enumerate}

\vspace{0.1in}

\boxguard[8]
\begin{spiralbox}
\small
\textbf{Next lesson (1.2, Change in Tandem):} instead of asking what the output \emph{is} at one
input, we ask what the output \emph{does} as the input climbs. Look back at your temperature
table in problem 4: when was the temperature rising, and when was it falling?
\end{spiralbox}

\end{document}
